\documentclass[11pt,a4paper]{article}

\usepackage[margin=2.6cm]{geometry}
\usepackage{amsmath,amssymb,amsthm,bm}
\usepackage{enumitem}
\usepackage{algorithm}
\usepackage{algpseudocode}
\usepackage{booktabs}
\usepackage{xcolor}
\usepackage[colorlinks=true,allcolors=blue]{hyperref}

\newtheorem{theorem}{Theorem}[section]
\newtheorem{lemma}[theorem]{Lemma}
\newtheorem{proposition}[theorem]{Proposition}
\newtheorem{corollary}[theorem]{Corollary}
\theoremstyle{definition}
\newtheorem{definition}[theorem]{Definition}
\newtheorem{remark}[theorem]{Remark}
\newtheorem{example}[theorem]{Example}

\newcommand{\R}{\mathbb{R}}
\newcommand{\eps}{\epsilon_{\mathrm{dbl}}}
\newcommand{\Mt}{\widetilde{\bm M}}
\newcommand{\normF}[1]{\|#1\|_{F}}
\newcommand{\rankn}{\operatorname{rank}_{\mathrm{num}}}

\title{Numerical Rank Certification of Interpolation Moment Matrices\\
by Column-Pivoted Householder QR Factorization}

\author{}
\date{}

\begin{document}
\maketitle
\vspace{-3em}

\begin{abstract}
The construction of generalized interpolation weights on scattered
three-dimensional stencils requires the solution of a dense
$27\times27$ moment system $\bm M\bm\Psi=\bm C$ with
$\bm M=\bm X^{\mathsf T}\bm W$, where $\bm X$ collects polynomial
samples up to degree six and $\bm W$ collects Hermite--Wendland radial
basis rows. Admissibility of a candidate stencil--parameter triple
$(\text{stencil},\beta,h)$ hinges on whether $\bm M$ possesses full
numerical rank in IEEE double precision. This note places the rank
test used in our solver on a rigorous and fully referenced footing. We
first review the Householder reflection and its role in orthogonal
triangularization (Section~\ref{sec:householder}), then the QR
factorization with column pivoting (QRCP) of Businger and Golub and
its rank-revealing interpretation (Section~\ref{sec:qrcp}), and
finally state the complete certification algorithm employed in our
code: iterative row/column max-norm equilibration, a three-stage
sequence of rank gates, and the acceptance criterion
$p_i>\tau_{\widetilde M}$ with
$\tau_{\widetilde M}=27\,\eps\,\normF{\Mt}$
(Section~\ref{sec:algorithm}). The threshold is shown to follow the
standard backward-error convention of numerical linear algebra, and
the result is characterized precisely as a double-precision numerical
certificate rather than a symbolic proof of nonsingularity.
\end{abstract}

%=====================================================================
\section{Householder reflections and orthogonal triangularization}
\label{sec:householder}

The elementary orthogonal transformation underlying all algorithms in
this note is the Householder reflection, introduced by
Householder~\cite{householder1958} as a tool for the unitary
triangularization of general matrices; systematic treatments are given
by Golub and Van Loan~\cite[\S5.1]{golubvanloan2013} and Trefethen and
Bau~\cite[Lecture~10]{trefethenbau1997}.

\begin{definition}[Householder reflector]
\label{def:householder}
Let $\bm v\in\R^{m}$, $\bm v\neq\bm 0$. The matrix
\begin{equation}
  \bm H \;=\; \bm I - 2\,\frac{\bm v\bm v^{\mathsf T}}
                              {\bm v^{\mathsf T}\bm v}
  \label{eq:H-def}
\end{equation}
is called a Householder reflector with Householder vector $\bm v$.
\end{definition}

Geometrically, $\bm H$ realizes the reflection of $\R^{m}$ across the
hyperplane $\{\bm x:\bm v^{\mathsf T}\bm x=0\}$ orthogonal to
$\bm v$~\cite[Lecture~10]{trefethenbau1997}. Two properties follow
directly from~\eqref{eq:H-def}: $\bm H$ is symmetric,
$\bm H=\bm H^{\mathsf T}$, and orthogonal,
$\bm H^{\mathsf T}\bm H=\bm I$, hence involutory,
$\bm H^{2}=\bm I$.

The utility of the reflector lies in its ability to annihilate all but
the first component of a prescribed vector.

\begin{lemma}[Annihilation property]
\label{lem:annihilation}
Let $\bm a\in\R^{m}$ with $\alpha=\|\bm a\|_{2}>0$ and
$\bm a\neq\alpha\,\bm e_{1}$. Choose
$\bm v=\bm a-\alpha\,\bm e_{1}$ and let $\bm H$ be the reflector
\eqref{eq:H-def}. Then
\begin{equation}
  \bm H\bm a \;=\; \alpha\,\bm e_{1}.
  \label{eq:Ha}
\end{equation}
\end{lemma}

\begin{proof}
Write $a_{1}=\bm e_{1}^{\mathsf T}\bm a$. Direct computation gives
\begin{equation}
  \bm v^{\mathsf T}\bm a
   =(\bm a-\alpha\bm e_{1})^{\mathsf T}\bm a
   =\alpha^{2}-\alpha a_{1},
\end{equation}
\begin{equation}
  \bm v^{\mathsf T}\bm v
   =(\bm a-\alpha\bm e_{1})^{\mathsf T}(\bm a-\alpha\bm e_{1})
   =\alpha^{2}-2\alpha a_{1}+\alpha^{2}
   =2\bigl(\alpha^{2}-\alpha a_{1}\bigr).
\end{equation}
Hence $\bm v^{\mathsf T}\bm a/\bm v^{\mathsf T}\bm v=\tfrac12$, and
therefore
\begin{equation}
  \bm H\bm a
   =\bm a-2\,\frac{\bm v^{\mathsf T}\bm a}{\bm v^{\mathsf T}\bm v}\,\bm v
   =\bm a-\bm v
   =\alpha\,\bm e_{1}. \qedhere
\end{equation}
\end{proof}

In finite-precision implementations the sign of $\alpha$ is chosen as
$\alpha=-\operatorname{sign}(a_{1})\|\bm a\|_{2}$ so that the
subtraction forming $\bm v$ incurs no cancellation; see
Higham~\cite[\S19.1]{higham2002} and
\cite[\S5.1.3]{golubvanloan2013}.

Applying Lemma~\ref{lem:annihilation} recursively to the trailing
submatrices of $\bm A\in\R^{m\times n}$ produces the Householder QR
factorization: after $n$ steps,
\begin{equation}
  \bm H_{n}\cdots\bm H_{2}\bm H_{1}\,\bm A=\bm R,
  \qquad
  \bm A=\bm Q\bm R,
  \qquad
  \bm Q=\bm H_{1}\bm H_{2}\cdots\bm H_{n},
\end{equation}
with $\bm R$ upper triangular and $\bm Q$ orthogonal. The decisive
advantage of this construction over Gram--Schmidt or Gaussian
elimination is its backward stability: the computed factors are the
exact factors of a perturbed matrix $\bm A+\Delta\bm A$ with
$\normF{\Delta\bm A}\le c(m,n)\,\eps\,\normF{\bm A}$ for a modest
polynomial $c(m,n)$, where $\eps\approx2.22\times10^{-16}$ denotes the
unit roundoff of IEEE binary64 arithmetic; see
Higham~\cite[Thm.~19.4]{higham2002} and
\cite[\S5.1.9]{golubvanloan2013}. This backward-error bound is the
quantitative basis of the rank threshold adopted in
Section~\ref{sec:algorithm}.

%=====================================================================
\section{QR factorization with column pivoting and numerical rank}
\label{sec:qrcp}

\subsection{The Businger--Golub algorithm}

The QR factorization with column pivoting (QRCP) was introduced by
Businger and Golub~\cite{busingergolub1965} in the context of linear
least-squares problems with (nearly) rank-deficient coefficient
matrices; the canonical modern description is
\cite[\S5.4.2]{golubvanloan2013}. Given $\bm A\in\R^{m\times n}$,
QRCP computes a permutation matrix $\bm P$, an orthogonal $\bm Q$, and
an upper-triangular $\bm R$ such that
\begin{equation}
  \bm A\bm P=\bm Q\bm R.
  \label{eq:qrcp}
\end{equation}
At step $k$ the algorithm selects, among the not-yet-processed
columns, the one whose trailing subcolumn (rows $k{:}m$) has largest
Euclidean norm, swaps it into position $k$, and annihilates its
subdiagonal entries with a Householder reflector
(Lemma~\ref{lem:annihilation}). The greedy pivot selection enforces
the monotonicity
\begin{equation}
  |r_{11}|\;\ge\;|r_{22}|\;\ge\;\cdots\;\ge\;|r_{qq}|,
  \qquad q=\min(m,n),
  \label{eq:monotone}
\end{equation}
and, more strongly, guarantees that each diagonal entry dominates the
trailing block column-wise,
$r_{kk}^{2}\ge\sum_{i=k}^{j}r_{ij}^{2}$ for all $j\ge
k$~\cite[\S5.4.2]{golubvanloan2013}. The reference implementation in
LAPACK is the routine \texttt{xGEQP3}~\cite{lapack1999}, the BLAS-3
formulation of Quintana-Ort\'{\i}, Sun, and
Bischof~\cite{quintanaorti1998}. The algorithm is deterministic: the
pivot at every step is defined by exact column-norm comparisons, so
repeated executions on the same data produce identical factorizations,
in contrast to randomized rank estimators.

\subsection{Mechanics of the pivot residual: trailing norms and their
evolution}
\label{sec:pivotresidual}

The quantity compared against the rank tolerance in
Section~\ref{sec:algorithm} is the \emph{pivot residual}
$p_k$ of step $k$. This subsection makes its definition, its
evolution under the reflections, and its geometric meaning precise,
since all three are essential to a correct reading of the acceptance
test: the residual is a property of the \emph{partially factored}
working matrix, not of the raw data.

Let $\bm A^{(0)}=\bm A\in\R^{m\times n}$ denote the (preprocessed)
input and let
\begin{equation}
  \bm A^{(k)}
  \;=\;
  \bm H_{k-1}\cdots\bm H_{1}\bm H_{0}\,\bm A\,\bm P_{0}\bm P_{1}\cdots
  \bm P_{k-1}
  \label{eq:workingmatrix}
\end{equation}
be the working matrix after $k$ pivoted reflection steps. For every
not-yet-selected column $j$, define its \emph{trailing norm} at step
$k$ as the Euclidean norm of its subvector below the already
triangularized block,
\begin{equation}
  \nu_k(j)
  \;=\;
  \Bigl(\textstyle\sum_{i=k}^{m-1}
        \bigl(a^{(k)}_{ij}\bigr)^{2}\Bigr)^{1/2},
  \qquad
  p_k=\max_{j\,\text{unselected}}\nu_k(j).
  \label{eq:trailingnorm}
\end{equation}
The leading components $a^{(k)}_{ij}$, $i<k$, are excluded: after the
first $k$ reflections they carry the coefficients of column $j$ with
respect to the already selected directions---entries of rows
$0,\dots,k-1$ of $\bm R$---and therefore contribute nothing new to
linear independence. At $k=0$ the trailing norm coincides with the raw
column norm; for every $k\ge1$ it is a quantity \emph{manufactured by
the preceding reflections}.

\begin{proposition}[Pivot residual equals the produced diagonal entry]
\label{prop:trailing}
At step $k$, let column $j^{\ast}$ attain the maximum in
\eqref{eq:trailingnorm} and be swapped into position $k$. Then the
diagonal entry created by the $k$th Householder reflection satisfies
\begin{equation}
  |r_{kk}|=\nu_k(j^{\ast})=p_k .
\end{equation}
Consequently the acceptance test $p_k>\tau$ may be evaluated
\emph{before} the reflection is formed, with a result identical to
inspecting $|r_{kk}|$ afterwards.
\end{proposition}

\begin{proof}
Write $\bm x\in\R^{m-k}$ for the trailing subvector of the pivot
column, so $\|\bm x\|_{2}=\nu_k(j^{\ast})$. By
Lemma~\ref{lem:annihilation} the reflector acting on components
$k,\dots,m-1$ maps $\bm x$ to $\alpha\,\bm e_{1}$ with
$|\alpha|=\|\bm x\|_{2}$, and $\alpha$ is precisely the entry stored
in position $(k,k)$ of $\bm R$. Orthogonality of the reflector
preserves the Euclidean length, so no other value is possible.
\end{proof}

Geometrically, the leading/trailing split of a working column is its
orthogonal decomposition with respect to the span of the selected
pivot columns:
\begin{equation}
  \underbrace{\bigl(a^{(k)}_{0j},\dots,a^{(k)}_{k-1,j}\bigr)}
             _{\text{component inside the selected subspace}}
  \;\Big|\;
  \underbrace{\bigl(a^{(k)}_{kj},\dots,a^{(k)}_{m-1,j}\bigr)}
             _{\text{component orthogonal to it}},
\end{equation}
so $\nu_k(j)$ is the Euclidean distance from column $j$ to the
subspace spanned by the $k$ columns already selected. A large
residual certifies a genuinely new direction; a residual below the
roundoff-level tolerance means the column lies, to double precision,
inside the subspace already captured---numerical linear dependence.
This is why the pivot residual measures \emph{independence} rather
than mere \emph{magnitude}, and why raw column norms cannot replace
it, as the following example shows.

\begin{example}[Raw norms cannot reveal rank]
\label{ex:nearparallel}
Let $\delta=10^{-15}$ and consider the two columns
\begin{equation}
  \bm a_{1}=(1,\,0,\,0)^{\mathsf T},
  \qquad
  \bm a_{2}=(1,\,\delta,\,0)^{\mathsf T},
\end{equation}
which are numerically parallel. Their raw norms are
$\|\bm a_{1}\|_{2}=1$ and $\|\bm a_{2}\|_{2}=\sqrt{1+\delta^{2}}
\approx1$; a test based on raw norms would count two directions and
report rank $2$. Under QRCP, step $0$ selects $\bm a_{1}$
($p_{0}=1$), and the reflection moves the component of $\bm a_{2}$
along $\bm a_{1}$ into the leading position, leaving the trailing
norm $\nu_{1}(2)=\delta=10^{-15}$. Against any tolerance of the
backward-error magnitude \eqref{eq:tolfamily} this residual fails,
$p_{1}\le\tau$, and the numerical rank is correctly reported as $1$.
The reflections thus convert magnitude information into independence
information; omitting them (or measuring full-length norms of the
unreflected data) invalidates the test from the second step onward.
\end{example}

Two further consequences of orthogonal invariance fix the roles of the
two norms appearing in the acceptance test $p_k>\tau$.

\begin{proposition}[The tolerance is frozen; the residuals evolve]
\label{prop:frozen}
The Frobenius norm is invariant under orthogonal transformations and
column permutations,
$\normF{\bm H\bm A\bm P}=\normF{\bm A}$, since
$\normF{\bm H\bm A}^{2}
=\operatorname{tr}(\bm A^{\mathsf T}\bm H^{\mathsf T}\bm H\bm A)
=\normF{\bm A}^{2}$. Hence the tolerance
$\tau=f(m,n)\,\eps\,\normF{\bm A^{(0)}}$, computed once from the
preprocessed input before the factorization begins, retains its value
in exact arithmetic throughout all steps; recomputing it would change
nothing except adding roundoff. The test $p_k>\tau$ therefore
compares an evolving measurement (the trailing norm, updated by every
reflection) against a fixed ruler (the roundoff-level noise bound of
the whole matrix).
\end{proposition}

\begin{remark}[Interleaved termination]
\label{rem:interleaved}
Textbook descriptions sometimes present rank determination as a
post-processing pass over the completed diagonal of $\bm R$. By
Proposition~\ref{prop:trailing} and the monotonicity
\eqref{eq:monotone}, the interleaved variant---test $p_k>\tau$ at
the top of every step and terminate at the first failure---produces
the identical rank count while offering two advantages. First,
efficiency: since $p_k$ is the \emph{maximum} remaining residual, its
failure guarantees that all subsequent pivots fail as well, so the
remaining reflections are provably useless work. Second, numerical
hygiene: forming a reflector from a pivot of roundoff-level norm
divides by the near-zero quantity $\bm v^{\mathsf T}\bm v$ and
sprays amplified noise into all remaining trailing blocks; stopping
before the degenerate reflection keeps the working matrix
uncontaminated. The rank-deficient least-squares driver
\texttt{xGELSY} of LAPACK adopts the same interleaved
strategy~\cite{lapack1999}.
\end{remark}

\subsection{Rank-revealing interpretation}

Because $\bm Q$ is orthogonal, the singular values of $\bm R$ coincide
with those of $\bm A$, and the profile of the diagonal
\eqref{eq:monotone} exposes the numerical rank: if
$|r_{kk}|$ is large for $k\le r$ and collapses to roundoff level for
$k>r$, then $\bm A$ has numerical rank $r$. Factorizations of the
form \eqref{eq:qrcp} whose triangular factor reliably separates the
dominant and negligible parts of the spectrum were named
\emph{rank-revealing QR} (RRQR) factorizations by
Chan~\cite{chan1987}. Interlacing arguments give the a priori bounds
\begin{equation}
  \sigma_{k}(\bm A)\;\ge\;|r_{kk}|
  \quad\text{is not guaranteed in general, but}\quad
  |r_{kk}|\;\le\;\sqrt{n-k+1}\,\sigma_{k}(\bm A),
\end{equation}
so a small pivot certifies a small singular value; the converse can
fail for contrived matrices. The classical counterexample is the
Kahan matrix~\cite{kahan1966}, for which the Businger--Golub pivots
decay far more slowly than the smallest singular value. Gu and
Eisenstat~\cite{gueisenstat1996} constructed \emph{strong} RRQR
algorithms with provable factor bounds that close this gap at modest
extra cost. In practice, however, QRCP fails to reveal rank only on
such adversarial examples; for equilibrated matrices of moderate
dimension it is the standard, inexpensive alternative to the singular
value decomposition, and it underlies the rank-deficient least-squares
driver \texttt{xGELSY} of LAPACK~\cite{lapack1999}.

\subsection{Numerical rank and tolerance conventions}

For a matrix stored and factored in floating-point arithmetic, rank is
meaningful only relative to a tolerance
$\tau$~\cite[\S5.4.1]{golubvanloan2013}: one defines
\begin{equation}
  \rankn(\bm A;\tau)
  \;=\;
  \#\{\,k:\ \sigma_{k}(\bm A)>\tau\,\}
  \quad\text{or, in the QRCP surrogate,}\quad
  \#\{\,k:\ p_{k}>\tau\,\},
  \label{eq:numrank}
\end{equation}
where $p_{k}$ denotes the $k$th pivot magnitude (the trailing
column-norm residual, equal to $|r_{kk}|$ at selection time). The
tolerance must dominate the backward error of the computation itself;
otherwise the test would attribute significance to quantities that are
indistinguishable from roundoff. Since Householder QR applied to
$\bm A\in\R^{n\times n}$ is backward stable with perturbation bound
proportional to $n\,\eps\,\|\bm A\|$~\cite[Thm.~19.4]{higham2002},
tolerances of the generic form
\begin{equation}
  \tau \;=\; f(m,n)\,\eps\,\|\bm A\|
  \label{eq:tolfamily}
\end{equation}
are the accepted convention, with $f$ a low-degree polynomial of the
dimensions and $\|\cdot\|$ a computable norm. Instances of this
family include the default tolerance
$\max(m,n)\,\eps\,\sigma_{\max}(\bm A)$ used for numerical rank
determination in standard software following
\cite[\S5.4.1]{golubvanloan2013}, and the relative-pivot condition
$|r_{kk}|/|r_{11}|>\mathrm{RCOND}$ of
\texttt{xGELSY}~\cite{lapack1999}. The particular member of the
family used in our solver is specified and justified next.

%=====================================================================
\section{The certification algorithm}
\label{sec:algorithm}

\subsection{Setting}

For every candidate triple $(\text{stencil},\beta,h)$ of a
$27$-point interpolation stencil, Wendland shape parameter
$\beta$, and support radius $h$, our generalized-interpolation
solver assembles a moment matrix
\begin{equation}
  \bm M=\bm X^{\mathsf T}\bm W\in\R^{27\times27},
\end{equation}
where the columns of $\bm X$ sample the polynomial basis of total
degree at most six on the stencil nodes and the columns of $\bm W$
contain Hermite--Wendland radial-basis
rows~\cite{wendland1995,wu1992}. The interpolation weights $\bm\Psi$
solve $\bm M\bm\Psi=\bm C$ for a prescribed moment vector $\bm C$. A
triple is \emph{admissible} if and only if this system is uniquely
solvable in double precision, i.e.\ if $\bm M$ has full numerical
rank $27$ in the sense of \eqref{eq:numrank}. Existence and
uniqueness of $\bm\Psi$ is governed by $\operatorname{rank}(\bm M)$
alone; the auxiliary tests on $\bm X$ and $\bm W$ described below are
diagnostic accelerators, not additional mathematical requirements.

\subsection{Row/column equilibration}
\label{sec:equilibration}

The entries of $\bm M$ mix polynomial degrees one through six and can
span many orders of magnitude. A pivot-magnitude test applied to the
raw matrix could then mistake benign scale imbalance for rank
deficiency, since column norms---the very quantities QRCP
pivots on---would reflect degree scaling rather than linear
independence. Diagonal scaling is the classical remedy: van der
Sluis~\cite{vandersluis1969} showed that norm equilibration brings the
condition number within a modest factor of its optimum over all
diagonal scalings. We employ the iterative row/column scheme in the
max norm, of the family analyzed by Ruiz~\cite{ruiz2001}: six sweeps,
each consisting of a row pass followed by a column pass,
\begin{equation}
\begin{aligned}
 s_i^{\mathrm{row}}&=\max_j|\widetilde M_{ij}|,&
 \widetilde M_{i:}&\leftarrow\widetilde M_{i:}/s_i^{\mathrm{row}},&
 d_i^{\mathrm{row}}&\leftarrow d_i^{\mathrm{row}}/s_i^{\mathrm{row}},\\
 s_j^{\mathrm{col}}&=\max_i|\widetilde M_{ij}|,&
 \widetilde M_{:j}&\leftarrow\widetilde M_{:j}/s_j^{\mathrm{col}},&
 d_j^{\mathrm{col}}&\leftarrow d_j^{\mathrm{col}}/s_j^{\mathrm{col}},
\end{aligned}
\label{eq:sweep}
\end{equation}
initialized with $\widetilde{\bm M}\leftarrow\bm M$ and
$d_i^{\mathrm{row}}=d_j^{\mathrm{col}}=1$. Only the diagonals are
stored, as vectors $\texttt{Dr}[27]$ and $\texttt{Dc}[27]$. After the
sweeps,
\begin{equation}
  \Mt=\bm D_r\bm M\bm D_c,
  \qquad
  \bm D_r=\operatorname{diag}(d_0^{\mathrm{row}},\dots,d_{26}^{\mathrm{row}}),
  \qquad
  \bm D_c=\operatorname{diag}(d_0^{\mathrm{col}},\dots,d_{26}^{\mathrm{col}}),
  \label{eq:equilibrated}
\end{equation}
with all diagonal entries positive and finite. Hence $\bm D_r$ and
$\bm D_c$ are nonsingular and
\begin{equation}
  \operatorname{rank}(\Mt)=\operatorname{rank}(\bm M),
\end{equation}
so equilibration alters the numerical scaling but not the exact rank.
Consistency is preserved when the linear system is eventually solved:
with $\widetilde{\bm C}=\bm D_r\bm C$, the solution $\bm y$ of
$\Mt\bm y=\widetilde{\bm C}$ yields $\bm\Psi=\bm D_c\bm y$, since
$\bm D_r\bm M\bm D_c\bm y=\bm D_r\bm C$ is equivalent to
$\bm M\bm\Psi=\bm C$.

\subsection{Three-stage rank gates and acceptance threshold}

Deterministic QRCP (Section~\ref{sec:qrcp}) is applied in three
stages. The first two stages test column-normalized copies of $\bm X$
and $\bm W$ individually; they serve as early-exit filters and as
diagnostics that attribute a failure either to degenerate stencil
geometry (rank loss in $\bm X$) or to an unsuitable
radial-basis parameter pair $(\beta,h)$ (rank loss in $\bm W$).
They are necessary conditions for the admissibility of the product,
because
\begin{equation}
  \operatorname{rank}(\bm X^{\mathsf T}\bm W)
  \;\le\;
  \min\{\operatorname{rank}(\bm X),\operatorname{rank}(\bm W)\},
\end{equation}
but they are not sufficient: full rank of both factors does not imply
full rank of their product. The decisive third stage therefore
applies QRCP to $\Mt^{\mathsf T}$ and counts as an independent
direction every pivot whose residual magnitude---the trailing norm
of Proposition~\ref{prop:trailing}, evaluated at the top of each step
and compared against the tolerance frozen at entry
(Proposition~\ref{prop:frozen}, Remark~\ref{rem:interleaved})---exceeds
\begin{equation}
  \boxed{\;
  p_i \;>\; \tau_{\widetilde M},
  \qquad
  \tau_{\widetilde M}\;=\;27\,\eps\,\normF{\Mt}.
  \;}
  \label{eq:threshold}
\end{equation}
The triple is accepted, and weights are exported, if and only if all
three gates report numerical rank $27$:
\begin{equation}
  \rankn(\bm X)=27,\qquad
  \rankn(\bm W)=27,\qquad
  \rankn(\Mt)=27.
\end{equation}
The complete procedure is summarized in
Algorithm~\ref{alg:certification}.

\begin{algorithm}[t]
\caption{Numerical rank certification of a candidate triple
$(\text{stencil},\beta,h)$}
\label{alg:certification}
\begin{algorithmic}[1]
\State Assemble $\bm X,\bm W\in\R^{27\times27}$;
       form $\bm M=\bm X^{\mathsf T}\bm W$
\State \textbf{Gate 1:} QRCP on column-normalized $\bm X$;
       \textbf{if} $\rankn(\bm X)<27$ \textbf{reject} (stencil geometry)
\State \textbf{Gate 2:} QRCP on column-normalized $\bm W$;
       \textbf{if} $\rankn(\bm W)<27$ \textbf{reject} (RBF parameters)
\State $\Mt\leftarrow\bm M$;\quad
       $\texttt{Dr}[i]\leftarrow1$, $\texttt{Dc}[j]\leftarrow1$
\For{$k=1,\dots,6$} \Comment{equilibration sweeps, Eq.~\eqref{eq:sweep}}
  \State row pass: divide each row of $\Mt$ by its max-norm;
         accumulate in \texttt{Dr}
  \State column pass: divide each column of $\Mt$ by its max-norm;
         accumulate in \texttt{Dc}
\EndFor
\State $\tau_{\widetilde M}\leftarrow27\,\eps\,\normF{\Mt}$
\State \textbf{Gate 3:} deterministic QRCP on $\Mt^{\mathsf T}$;
       count pivots with $p_i>\tau_{\widetilde M}$
\If{count $=27$}
  \State \textbf{accept}; on demand solve
         $\Mt\bm y=\bm D_r\bm C$ and return
         $\bm\Psi=\bm D_c\bm y$
\Else
  \State \textbf{reject}; advance to the next
         $(\text{stencil},\beta,h)$ trial
\EndIf
\end{algorithmic}
\end{algorithm}

\subsection{Justification of the threshold}

The tolerance \eqref{eq:threshold} is the member of the standard
family \eqref{eq:tolfamily} obtained with dimension factor
$f=n=27$ and the Frobenius norm as the computable measure of
$\|\Mt\|$. Its role follows directly from backward stability: the
computed factorization corresponds exactly to a perturbed matrix
$\Mt+\Delta$ with $\normF{\Delta}\le c(n)\,\eps\,\normF{\Mt}$
\cite[Thm.~19.4]{higham2002}, so any pivot residual below the level
$n\,\eps\,\normF{\Mt}$ lies within the uncertainty introduced by the
arithmetic itself and cannot be distinguished from a numerically zero
direction. Two remarks make the choice precise.

\begin{remark}[Conservativeness of the Frobenius norm]
Since $\|\bm A\|_{2}\le\normF{\bm A}\le\sqrt{n}\,\|\bm A\|_{2}$, the
tolerance $27\,\eps\,\normF{\Mt}$ is at least as large as the
spectral-norm variant $27\,\eps\,\sigma_{\max}(\Mt)$
recommended in \cite[\S5.4.1]{golubvanloan2013}. The test is
therefore conservative---biased toward rejection---while avoiding
the cost of estimating $\sigma_{\max}$. After the equilibration of
Section~\ref{sec:equilibration} the entries of $\Mt$ are of order
unity, so the two norms differ by at most the benign factor
$\sqrt{27}$.
\end{remark}

\begin{remark}[Nature of the certificate]
\label{rem:certificate}
A passed Gate~3 asserts that all $27$ QRCP pivots of $\Mt^{\mathsf T}$
exceed the roundoff-level tolerance \eqref{eq:threshold} in IEEE
double precision. In view of the Kahan
counterexample~\cite{kahan1966} and the strong-RRQR theory of Gu and
Eisenstat~\cite{gueisenstat1996}, this constitutes a
\emph{double-precision numerical certificate} of
$\rankn(\bm M)=27$---sufficient for the reliable solution of
$\bm M\bm\Psi=\bm C$ in the same arithmetic---and not a symbolic proof
of $\det(\bm M)\neq0$. For equilibrated matrices of dimension $27$
arising from non-adversarial stencil geometries, the gap between the
two notions is immaterial in practice, and QRCP-based rank
determination is the established inexpensive surrogate for the
singular value decomposition~\cite{chan1987,lapack1999}.
\end{remark}

%=====================================================================
\section*{Concluding note}

The certification procedure of Section~\ref{sec:algorithm} is built
entirely from classical, well-audited components: Householder
reflections~\cite{householder1958,golubvanloan2013,trefethenbau1997},
the Businger--Golub column-pivoted QR
factorization~\cite{busingergolub1965,quintanaorti1998,lapack1999},
rank-revealing interpretation and its known
limitations~\cite{chan1987,kahan1966,gueisenstat1996}, iterative
max-norm equilibration~\cite{vandersluis1969,ruiz2001}, and the
standard backward-error tolerance convention~\cite{higham2002,
golubvanloan2013}. The only problem-specific ingredients are the
dimension $n=27$, fixed by the D3Q27 interpolation stencil, and the
three-gate ordering, which converts the mathematical
characterization---unique solvability if and only if
$\operatorname{rank}(\bm M)=27$---into an efficient and diagnosable
selection loop over candidate triples $(\text{stencil},\beta,h)$.

%=====================================================================
\begin{thebibliography}{99}

\bibitem{householder1958}
A.~S. Householder,
\newblock Unitary triangularization of a nonsymmetric matrix,
\newblock \emph{Journal of the ACM}, 5(4):339--342, 1958.

\bibitem{busingergolub1965}
P.~Businger and G.~H. Golub,
\newblock Linear least squares solutions by {H}ouseholder
transformations,
\newblock \emph{Numerische Mathematik}, 7:269--276, 1965.

\bibitem{golubvanloan2013}
G.~H. Golub and C.~F. Van~Loan,
\newblock \emph{Matrix Computations},
\newblock 4th edition, Johns Hopkins University Press, Baltimore,
2013. (QR with column pivoting: \S5.4.2; numerical rank and the SVD:
\S5.4.1; Householder reflections: \S5.1.)

\bibitem{trefethenbau1997}
L.~N. Trefethen and D.~Bau, III,
\newblock \emph{Numerical Linear Algebra},
\newblock SIAM, Philadelphia, 1997. (Householder triangularization:
Lecture~10.)

\bibitem{higham2002}
N.~J. Higham,
\newblock \emph{Accuracy and Stability of Numerical Algorithms},
\newblock 2nd edition, SIAM, Philadelphia, 2002. (Backward stability
of Householder QR: Chapter~19, Theorem~19.4.)

\bibitem{chan1987}
T.~F. Chan,
\newblock Rank revealing {QR} factorizations,
\newblock \emph{Linear Algebra and its Applications},
88--89:67--82, 1987.

\bibitem{gueisenstat1996}
M.~Gu and S.~C. Eisenstat,
\newblock Efficient algorithms for computing a strong rank-revealing
{QR} factorization,
\newblock \emph{SIAM Journal on Scientific Computing},
17(4):848--869, 1996.

\bibitem{kahan1966}
W.~Kahan,
\newblock Numerical linear algebra,
\newblock \emph{Canadian Mathematical Bulletin}, 9:757--801, 1966.

\bibitem{quintanaorti1998}
G.~Quintana-Ort\'{\i}, X.~Sun, and C.~H. Bischof,
\newblock A {BLAS}-3 version of the {QR} factorization with column
pivoting,
\newblock \emph{SIAM Journal on Scientific Computing},
19(5):1486--1494, 1998.

\bibitem{lapack1999}
E.~Anderson, Z.~Bai, C.~Bischof, S.~Blackford, J.~Demmel, J.~Dongarra,
J.~Du~Croz, A.~Greenbaum, S.~Hammarling, A.~McKenney, and D.~Sorensen,
\newblock \emph{{LAPACK} Users' Guide},
\newblock 3rd edition, SIAM, Philadelphia, 1999. (Routines
\texttt{xGEQP3}, \texttt{xGELSY}.)

\bibitem{vandersluis1969}
A.~van~der~Sluis,
\newblock Condition numbers and equilibration of matrices,
\newblock \emph{Numerische Mathematik}, 14:14--23, 1969.

\bibitem{ruiz2001}
D.~Ruiz,
\newblock A scaling algorithm to equilibrate both rows and columns
norms in matrices,
\newblock Technical Report RAL-TR-2001-034, Rutherford Appleton
Laboratory, Oxfordshire, 2001.

\bibitem{wendland1995}
H.~Wendland,
\newblock Piecewise polynomial, positive definite and compactly
supported radial functions of minimal degree,
\newblock \emph{Advances in Computational Mathematics}, 4:389--396,
1995.

\bibitem{wu1992}
Z.~Wu,
\newblock Hermite--{B}irkhoff interpolation of scattered data by
radial basis functions,
\newblock \emph{Approximation Theory and its Applications},
8(2):1--10, 1992.

\end{thebibliography}

\end{document}
